\documentclass{article}

% *** ACTION NEEDED BEFORE SUBMISSION ***
% The Meta-Agents CFP requires DOUBLE-blind review (authors, affiliations, and
% acknowledgments removed; prior work cited in third person). The style option
% below ("sglblindworkshop") is SINGLE-blind and must be swapped for the
% workshop template's genuine double-blind/anonymous option -- check the
% option name in the .sty file you download from the workshop site, since it
% is not verified here.
% natbib is loaded by the style file; options must be passed ahead of it.
\PassOptionsToPackage{numbers,sort&compress}{natbib}
\usepackage[dblblindworkshop]{neurips_2026} % TODO: replace with double-blind option
% \workshoptitle{...} % surfaces in the footer only once the [final] option is added

\usepackage[utf8]{inputenc}
\usepackage[T1]{fontenc}
\usepackage{microtype}

% Math & symbols
\usepackage{amsmath,amssymb}

% Graphics, tables, code
\usepackage{graphicx}
\usepackage{booktabs}
\usepackage{array}
\usepackage{tabularx}
\usepackage{listings}
\usepackage{xcolor}

% Diagrams
\usepackage{tikz}
\usetikzlibrary{positioning,fit,backgrounds,arrows.meta,calc,shadows}

% Links (loaded after the style file; cleveref must come last)
\usepackage[hidelinks]{hyperref}
\usepackage{cleveref}

% Personal comment commands removed for double-blind submission (do not
% reintroduce named-author comments in the submitted version).
% TODO: delete before submission
\newcommand{\SC}[1]{\textcolor{violet}{Simon: #1}}
\newcommand{\JB}[1]{\textcolor{teal}{Jaloliddin: #1}}

% Convenience macros (mirrors main.tex)
\usepackage{xspace}
\newcommand{\rlm}{\textsc{RLM}\xspace}
\newcommand{\rlms}{\textsc{RLM}s\xspace}
\newcommand{\sh}{\textsc{Self-Harness}\xspace}
\newcommand{\sid}[1]{\textbf{\texttt{#1}}}

% ---- styling for the RLM edit-space figure (fig:surfaces) ----
\definecolor{editgreen}{HTML}{2E7D5B}
\definecolor{editfill}{HTML}{E7F5EE}
\definecolor{fixgrey}{HTML}{8A8F98}
\definecolor{fixfill}{HTML}{F1F2F4}
\definecolor{envtint}{HTML}{EAF1FB}
\definecolor{subtint}{HTML}{FCF1E7}
\definecolor{inkcol}{HTML}{2B2F36}
\newcommand{\cardhd}{\bfseries}
\newcommand{\cardsb}{\normalfont\scriptsize\color{black!62}}
\tikzset{
  penbadge/.pic={
    \fill[editgreen] (0,0) circle (0.095);
    \draw[white, line width=0.55pt, line cap=round] (-0.045,-0.045) -- (0.03,0.03);
    \fill[white] (-0.052,-0.052) circle (0.017);
  },
  lockbadge/.pic={
    \fill[fixgrey] (0,0) circle (0.095);
    \draw[white, line width=0.6pt] (-0.032,0.006) arc (180:0:0.032);
    \fill[white, rounded corners=0.2pt] (-0.05,-0.05) rectangle (0.05,0.014);
  },
}

% The initial harness: the sparse floor the optimization loop starts from.
\newcommand{\hzero}{\ensuremath{\mathrm{H}_{0}}\xspace}
\newcommand{\hzerostar}{\ensuremath{\mathrm{H}_{0}^{*}}\xspace}
% Figure_1.tex is a tall float; let it take more of a text page so it lands near
% the discussion instead of being deferred to a page of its own.
\renewcommand{\topfraction}{0.9}
\renewcommand{\textfraction}{0.07}
\renewcommand{\floatpagefraction}{0.85}

% ---- code-listing style for the initial-harness figure (Figure_1.tex) ----
\definecolor{lstbg}{HTML}{FAFAFB}
\definecolor{lstframe}{HTML}{D5D8DD}
\definecolor{lstkw}{HTML}{1052B8}
\definecolor{lststr}{HTML}{B25A05}
\definecolor{lstnum}{HTML}{9AA0A8}
\definecolor{lstcom}{HTML}{6B7280}
\lstdefinestyle{harness}{
  language=Python,
  basicstyle=\ttfamily\scriptsize\linespread{0.88}\selectfont,
  keywordstyle=\color{lstkw},
  stringstyle=\color{lststr},
  commentstyle=\color{lstcom},
  deletekeywords={None,True,False},
  emph={str,list,dict,int,tuple},
  emphstyle=\color{lstkw},
  numbers=left,
  numberstyle=\ttfamily\tiny\color{lstnum},
  numbersep=6pt,
  backgroundcolor=\color{lstbg},
  frame=single,
  rulecolor=\color{lstframe},
  framesep=3pt,
  breaklines=true,
  breakindent=14pt,
  breakautoindent=false,
  columns=fullflexible,
  keepspaces=true,
  showstringspaces=false,
  xleftmargin=15pt,
  xrightmargin=1pt,
  aboveskip=1pt,
  belowskip=1pt,
}

\title{Self-Harnessing Recursive Language Models}
% \author{%
%   William Stanford \\
%     Affiliation  \\
%     Address \\
%     \texttt{email} \\
%     \And
%     Jaloliddin Boymakhammadov \\
%     Affiliation \\
%     Address \\
%     \texttt{email} \\
%     \AND
%     Eliaz Calvar \\
%     Affiliation \\
%     Address \\
%     \texttt{email} \\
%     \And
%     Mohammed Akram \\
%     Affiliation \\
%     Address \\
%     \texttt{email} \\
%     \And
%     Simon Coumes \\
%     Affiliation \\
%     Address \\
%     \texttt{email} \\
% }

\begin{document}

\maketitle

\begin{abstract}
Recursive language models (\rlms) let a bounded-context model operate over inputs larger than
its context window by decomposing it, invoking recursive model calls, and aggregating their
outputs~\citep{zhang2025rlm}, but vanilla RLMs fail in recurring ways that so far require
fine-tuning the model or hand-extending the harness. We propose adapting the \sh{} framework
of \citet{zhang2026selfharness} to \rlms, so a fixed model improves its own recursive harness
from evidence produced while running: it mines recurring failure mechanisms from
verifier-grounded execution traces, proposes minimal edits to declared harness surfaces, and
promotes only edits that survive non-regressive validation. The optimized harness is frozen
and evaluated on instances $8$--$32\times$ longer than any seen during optimization and in new
target environments. Results forthcoming.
\end{abstract}

\section{Introduction}
Frontier language models are bounded by a fixed context window, and quality degrades on
inputs longer than what the model was effectively trained and evaluated on, well before the
window limit is reached, a failure mode termed \emph{context rot}~\citep{hong2025contextrot}.
An input is called \emph{in-distribution} when its length and structure resemble what the model reliably handles, and \emph{out-of-distribution} otherwise, where behavior degrades unpredictably. Recursive language models (\rlms) scale inference-time compute instead of the
context window~\citep{snell2024scaling}: a model decomposes its own context and invokes fresh
copies of itself on the pieces~\citep{zhang2025rlm}, turning a long out-of-distribution input
into sub-calls that are individually shorter~\citep{zhang2026harness}; shortening addresses
the length dimension of the mismatch, though it does not by itself guarantee a sub-call's
content is otherwise in-distribution.
Performance therefore depends on the harness governing recursion, not just the model. 
Policies collapse into a single sub-call~\citep{zhang2026harness},
deeper recursion degrades accuracy while inflating cost~\citep{wang2026overthink}, and recursion
halts by heuristic rather than evidence~\citep{guha2026termination}, failures that implicate
the harness.

These three approaches are mainly used today.
Fine-tuning the model to work within the
harness~\citep{zhang2026harness,yang2026recursive} ties any gain to a specific checkpoint.
Hand-engineering the harness~\citep{alizadeh2026srlm,roy2026ycombinator,lumer2026rah} requires
frequent expert work that a new model family can outpace. 
A third, \emph{meta-agent} approach edits the scaffold directly by inspecting an agent's own trajectories, as in \sh{}~\citep{zhang2026selfharness}, which mines execution traces for recurring failures, proposes minimal edits, and promotes only non-regressive ones. We adapt \sh{} to \rlms below. \rlms are a natural but untested setting for this third approach: failures are individually checkable, grounding attribution in the call tree rather than a proposer's reading of a trace. 
Prime Agent's Continual Harness lets an RLM accumulate deployment state (history, memory, skills, subagent specs) around a fixed control loop, rather than editing the loop's decision logic directly~\citep{karten2026primeagent}.
Agentic Harness Engineering is the closest
instantiation of the propose-evidence-promote pattern we adapt to RLMs, showing a frozen
evolved harness transfers across benchmarks and model families for coding
agents~\citep{lin2026agenticharness}. Neither has been tested where recursion depth and
sub-call decomposition are themselves modifiable by the meta agent. 
We ask: \emph{can a language model
identify recurring failures in its own recursive behavior, convert them into targeted harness
edits, and build a harness that keeps performing when applied to inputs longer than those it was optimized on?}

\section{Our Contributions}

\begin{enumerate}
\setlength{\itemsep}{1pt}
\setlength{\topsep}{2pt}
\setlength{\parskip}{0pt}
\item \textbf{\sh{} for recursive language models.}
Adapting \sh{} \citep{zhang2026selfharness} to \rlms by declaring ten editable harness surfaces keyed to phases of the \rlm{} turn loop, starting from a deliberately sparse initial harness and letting a fixed model mine its own traces, propose targeted edits, and validate them without weight updates or a stronger external model.

\item \textbf{Evidence of length generalization.}
Optimising only on short source instances, freeze the harness, and evaluate on source instances $8$--$32\times$ longer.

\item \textbf{Evidence of cross-environment transfer.}
Applying the same frozen harness to an unseen target environment's short and long splits.
\end{enumerate}

\section{Preliminaries and Self-Harness Optimization}
\label{sec:method}

An \rlm{} answers a query by driving a Python REPL (read-eval-print loop) turn by turn: on each
turn the model writes one code block, sees only that block's printed output, and works toward a
final answer accumulated in a REPL variable~\citep{zhang2025rlm}. We call this sequence, from
query to submitted answer, the \rlm{} \textbf{turn loop}. Three invariants of the reference
implementation stay fixed throughout: \textbf{prompt-as-variable} (the prompt lives in the REPL
as a variable, never copied into the root context), \textbf{programmatic sub-calls} (recursive
calls to fresh copies of the model are issued by code, in loops, over slices of the prompt
variable), and \textbf{outputs-in-variables} (the final answer accumulates in, and is returned
from, a REPL variable). \sh{} edits the policies governing this loop, not the model's weights: it acts through four runtime injection points, an injected metadata function, answer-protocol middleware, sub-call retry/validation, and enforceable batch caps, each of which defaults to the reference implementation's shipped behavior. These hooks let SELF-HARNESS tune the loop's behavior without altering the three invariants above, the model's weights, external tools, or the evaluator, all of which remain fixed throughout.


Following \citet{zhang2026selfharness}, we declare \textbf{ten editable surfaces} \cref{app:surfaces}, one for each
phase of the turn loop above. Each surface is implemented as
a single builder function that \sh{} may rewrite. The surfaces are declared
directly from the loop's phase structure rather than selected from a published catalogue of
known \rlm{} failure modes, so anything the optimization loop discovers can be attributed to
the loop itself rather than to a designer who already knew what to fix. Optimization is started
from an initial harness \hzero{}:
seven of the ten surfaces return nothing, a disabled policy, or a single generic line, so any
orchestration guidance a tuned harness carries must be recovered from the loop's own execution
traces rather than inherited from \hzero{}.


\subsection{Self-Harness Optimization}
\label{sec:sho}

SH-RLMs adapt the \sh{} framework of \citet{zhang2026selfharness} to \rlms: the same fixed
model executes tasks and proposes harness changes. A proposal modifies exactly one of the ten
declared surfaces. One further surface is permanently excluded: letting the proposer truncate or
pre-summarize the prompt variable into root context, which would violate the first invariant by
turning the \rlm{} back into a compaction agent. Each round runs the three following stages.

\textbf{Weakness mining} runs the current harness on short mining-data instances, recording the verifier outcome and full recursive trace of every run. Each sub-call is scored by the environment's sub-verifier; since sub-problems are authored by the model itself, a sub-call may be correct, incorrect, or uncheckable. Mining is additionally run with the sub-verifier withheld, so the proposer sees only root outcomes and raw traces; execution is otherwise identical, so the two mining passes differ only in whether child-level evidence is available to the same fixed model judging causal status. Each failure becomes a structured record: the verifier-level cause and failing level come from the verifier and sub-verifier, while the causal status and implicated mechanism are judged by the same fixed model from a compressed trace digest, validated against a closed mechanism vocabulary mapped one-to-one to editable surfaces. Failures are clustered by exact agreement on the verifier-grounded signature $\phi(r_i)=(\textit{verifier cause}, \textit{failing level}, \textit{causal status}, \textit{agent mechanism})$, and the resulting evidence bundle ranks each recurring pattern by instances affected and estimated actionability, describing weaknesses without prescribing edits.

\textbf{Harness proposal} gives the same fixed model the failure patterns of the evidence bundle, a record of passing behaviors to preserve, and a summary of previously attempted edits, and asks it to generate several distinct, minimal candidate edits. Each proposal must target one mined failure pattern, modify one declared surface, state its predicted effect, and identify possible regressions.


\textbf{Proposal validation} evaluates the incumbent harness and each candidate on the mining data and a disjoint validation split never shown to the proposer, running a preregistered number of repetitions per instance. Candidates that fail a structural check (a stale base harness, zero or multiple surfaces changed, an invariant violation) are rejected without evaluation. For every remaining candidate we compute two numbers relative to the incumbent: $\Delta_{\text{mine}}$, the change in pass count on the mining data, and $\Delta_{\text{val}}$, the change in pass count on the validation data. A candidate is promoted only if it satisfies three conditions: it does not regress on either split by more than a preregistered tolerance $\tau_{\text{reg}}$ (that is, $\Delta_{\text{mine}} \geq -\tau_{\text{reg}}$ and $\Delta_{\text{val}} \geq -\tau_{\text{reg}}$); it improves on at least one split by more than a second preregistered threshold $\tau_{\text{imp}}$ (that is, $\max(\Delta_{\text{mine}}, \Delta_{\text{val}}) > \tau_{\text{imp}}$); and its mean sub-call count and cost fall within a preregistered band. Both tolerances are fixed from pilot variance before optimization begins, so they cannot be tuned to favor a particular candidate; setting them to zero recovers the strict no-regression rule of Zhang et al.~\cite{zhang2026selfharness}. Evaluation limits are enforced outside the harness, so a candidate's runtime policy may tighten these limits but never exceed them. Unlike Zhang et al.~\cite{zhang2026selfharness}, when multiple accepted edits touch disjoint surfaces we re-evaluate the merged harness before promotion; if the merge fails, the round promotes nothing. Optimization stops after a fixed number of rounds or several consecutive rounds without promotion.


\section{Experiments}

\subsection{Setup}
\label{sec:setup}
The experiments are a two-stage procedure: (1) a \sh{} optimization stage, in which a frozen-weight RLM edits its own harness against mining and validation data, and (2) a fixed-weight evaluation stage, comparing the frozen harness against fixed baselines on held-out data never touched during optimization. All conditions share a single frozen backbone, \textbf{Deepseek V4 Flash}, with the same configuration for root and sub-calls; only the harness code differs between conditions.

\subsection{Environments and splits}The setup is evaluated on two environments from the length-generalization suite of \citet{zhang2026harness}. \textbf{GraphWalks}~\citep{openai2025graphwalks} gives a graph and asks for all nodes reachable from a source via multi-hop traversal; a sub-call restricted to an induced subgraph is independently checkable. \textbf{OOLONG-Pairs} adapts OOLONG~\citep{bertsch2025oolong}, a long-context benchmark built from many short, independent records, by asking which pairs of records jointly satisfy a stated relational constraint; each candidate pair is checkable in isolation. Both tasks thus decompose into independently solvable units and admit \emph{synthesized sub-verifiers} that check a sub-problem without the root answer. Both also provide matched short/long instances ($8$--$32\times$ larger) under an unmodified deterministic verifier: GraphWalks is the source (optimization) environment, OOLONG-Pairs the target (held-out, evaluation-only).

Only the source environment's short split is used for optimization, divided into \textbf{mining data} ($n_{\text{in}}=24$, used in weakness mining and proposal validation) and \textbf{validation data} ($n_{\text{ho}}=40$, used only during proposal validation). Everything else, the remaining source-short test instances, all of source-long and both splits of the target environment is \textbf{held-out data}: untouched during optimization and used only for the final evaluation, across four categories, each with 40 short / 150 long test instances: \textbf{source-short} (optimization-length improvement), \textbf{source-long} (length generalization), \textbf{target-short} (cross-environment transfer), and \textbf{target-long} (transfer under an added length shift).

\input{figure_short_paper}
\subsection{Baselines and Evaluation}
\label{sec:baselines}

SH-RLM is evaluated against three fixed-weight baselines: \textbf{B1, initial \rlm{} (\hzero{})}, the unmodified reference harness and the starting point of optimization; \textbf{H0*, \rlm{} reference}~\citep{zhang2025rlm}, the hand-designed upstream \rlm{} harness, representing expert prompt and orchestration engineering within the same free-form paradigm as \hzero{}; and \textbf{$\lambda$-RLM~\citep{roy2026ycombinator}}, a hand-designed extension replacing free-form recursive code generation with a typed functional runtime. Together, H0* and $\lambda$-RLM measure \sh{} against two different styles of human harness engineering, a hand-tuned prompting within the same execution paradigm, and a structurally different typed alternative.
% rather than only against our own sparse starting point. 
% One further condition relaxes the fixed-weight constraint: \textbf{F1, fine-tuned \rlm{}}, the same backbone RL-trained inside the initial harness on the source short split (\cref{app:finetune}), bounding how much of the residual gap owes to model capability rather than orchestration; F1 is excluded from the primary comparisons.

The primary comparisons are SH-RLM vs.\ B1, vs.\ H0*, and vs.\ $\lambda$-RLM on source-long, target-short, and target-long: B1 measures improvement over the initial harness, while H0* and $\lambda$-RLM test whether automatic harness optimization can compete with manual harness and algorithm design; source-long tests whether learned changes survive a length shift, target-short is the cleanest evidence of transfer, and target-long tests transfer under simultaneous shifts.

The primary metric is verifier accuracy on all four held-out test categories, averaged over seeded runs with bootstrap confidence intervals; SH-RLM is compared against each baseline. Secondary metrics are token counts, recursive-call count, recursion depth, and accuracy per million tokens. Trace analysis reports mined-failure-pattern frequency before/after optimization, whole-input sub-call collapse, and (via the sub-verifiers) the root/child share of failures, showing whether optimization repaired errors or relocated them.

\paragraph{Status.} Optimization and evaluation infrastructure for both environments is implemented and weakness mining is operational; harness proposal, validation, and the four-way evaluation are in progress. Results are pending completion of the experiments described above.

\section{Conclusion and Future Work}

SH-RLM tests whether a fixed model can mine its own recursive execution traces, propose
minimal edits to a declared set of harness surfaces, and promote only validated ones, then
generalize the frozen result across a length shift and an unseen environment. A positive
result would show the propose-validate-promote loop transfers beyond the flat, tool-based
action spaces where it has so far been demonstrated to a harness whose editable surfaces are
tied to a recursive control loop, and is competitive with both styles of hand-designed harness engineering we test against the same-paradigm reference implementation (H0*) and the structurally different typed alternative, $\lambda$-RLM without any weight update. We leave the fine-tuned RLM baseline (F1, Appendix C) to future work: it would test whether SH-RLM's gains come mainly from editing the harness, or whether training the same starting model further would close a similar share of the gap. 

\paragraph{Responsible-use statement.}
SH-RLM can be used in the future as an add-on for RLMs: once results are available, it may
help improve the functionality of standard RLMs on long-context tasks. All experiments use
only synthetic or public benchmark data, with no personal or sensitive information, and no
harness edit is deployed or used to make decisions affecting real people; edits are
evaluated solely against benchmark verifiers. Given this scope, we do not see a direct
societal risk from the work at this stage; \cref{app:alignment} discusses a safety-monitoring
extension relevant should self-editing harnesses of this kind be deployed in the future. 

\clearpage
\bibliographystyle{unsrtnat}
\bibliography{references}

\clearpage
\appendix

\section{Editable Surface Reference}
\label{app:surfaces}

\begin{table}[htbp]
    \centering
    \renewcommand{\arraystretch}{1.4}
    \begin{tabular}{p{0.25\textwidth} p{0.70\textwidth}}
        \hline
        \textbf{Surface} & \textbf{Governs} \\
        \hline
        \textbf{S1} --- REPL contract & The factual contract given to the root model: available names, the answer protocol, per-turn REPL-block and stdout conventions, truncation behavior. \\
        \hline
        \textbf{S2} --- Decomposition instruction & Turns 1--2: how the root probes the input and whether/how it plans a decomposition. \\
        \hline
        \textbf{S3} --- Execution instruction & Per-turn discipline: what to print, when to offload to a sub-call vs. read directly, how to aggregate results. \\
        \hline
        \textbf{S4} --- Verification instruction & What the root checks before marking an answer ready. \\
        \hline
        \textbf{S5} --- Recovery instruction & What the root does when a sub-call errors or returns something unusable. \\
        \hline
        \textbf{S6} --- Runtime policy & Numeric limits and switches: chars per prompt, batch width, max calls per turn/total, recursion depth, retry-on-error, sub-output validation. \\
        \hline
        \textbf{S7} --- Metadata function & What carries across turns --- the harness's memory of prior calls. \\
        \hline
        \textbf{S8} --- REPL helpers & Harness-local functions and data injected into the REPL namespace (chunkers, batch wrappers, etc.). \\
        \hline
        \textbf{S9} --- Answer  middle-ware & Programmatic inspection of the detected final answer, with the ability to redirect (suppress and continue) rather than accept it. \\
        \hline
        \textbf{S10} --- Skills &
        Named, reusable procedures available across turns. The system prompt exposes
        a name--description index, while full skill bodies are loaded on demand through
        \texttt{load\_skill(name)} in the root and recursive child REPLs. \\
        \hline
    \end{tabular}
    \caption{Overview of the ten editable RLM harness surfaces and the behavior governed by each.}
    \label{tab:surfaces}
\end{table}

% Figure 1: the initial harness H0 and its ten editable surfaces.
% Body only; the `harness' listing style, its colours, and \hzero live in the
% preamble of proposal.tex, matching how Figure_2.tex relies on the tikz styling
% defined there.
\begin{figure}[!tbp]
    \centering
    \begin{lstlisting}[style=harness]
    # shrlm/rlm_harness.py: the ten surfaces Self-Harness may edit (S1 to S10).
    
    # Condensed for visibility; the full literal appears in the appendix.
    MINIMAL_REPL_CONTRACT = """
    You are an RLM operating through a persistent Python REPL.
    Use `context`, flat or recursive LM calls, `SHOW_VARS()`, and `answer`.
    Execute one ```repl``` block per turn; print inspections and mark the
    answer ready only when complete. [Display-only condensation.]
    """
    
    def build_repl_contract() -> str:                                    # S1
        return MINIMAL_REPL_CONTRACT
    
    def build_decomposition_instruction() -> str:                        # S2
        return ("Start by probing `context` to understand what you have, "
                "then decide how the task breaks into steps.")
    
    def build_execution_instruction() -> str:                            # S3
        return ("On each turn, run one block of code, print what you need "
                "to see, and build the result up in REPL variables.")
    
    def build_verification_instruction() -> str:                         # S4
        return 'Check your candidate answer before setting `answer["ready"] = True`.'
    
    def build_recovery_instruction() -> str:                             # S5
        return ("If a sub-call errors or returns something you cannot use, "
                "adjust the approach and try again.")
    
    def build_runtime_policy() -> dict[str, Any]:                        # S6
        return {
            "enabled": False,
            "max_prompt_chars": None,
            "max_batch_width": None,
            "max_depth": None,
            "retry_on_syntax_error": None,
            "max_retries": None,
            "validate_sub_output": None,
        }
    
    def build_metadata(                                                  # S7
        stdout: str,
        repl_inventory: dict[str, tuple[str, int]],
        max_character_length: int = DEFAULT_MAX_CHARACTER_LENGTH,
    ) -> str:
        return default_metadata_builder(stdout, repl_inventory,
                                        max_character_length)
    build_metadata.declared_bound = DEFAULT_MAX_CHARACTER_LENGTH        # S7
    
    def build_repl_helpers() -> dict[str, Any]:                          # S8
        return {}
    
    def build_sub_repl_helpers() -> dict[str, Any]:                      # S8
        return {}
    
    def build_answer_middleware() -> AnswerMiddlewareFn:                 # S9
        return accept_answer
    
    @dataclass(frozen=True)
    class SkillEntry:                                                    # S10
        name: str          # REPL-safe identifier; index label
        description: str   # one line: when to consult it; enters the prompt
        body: str          # ordered steps; served verbatim by `load_skill(name)`
    
    def build_skills() -> list[SkillEntry]:                              # S10
        return []
    \end{lstlisting}
    \caption{Initial harness \hzero{}, the starting point for \sh{}: ten declared
    surfaces, one builder function each, keyed to phases of the \rlm{} turn loop.
  %  Eight of the ten return nothing, a disabled policy, or a single generic line, so the orchestration strategy a tuned harness carries must be recovered by the optimization loop from its own execution traces. 
    The \sid{S1} contract is
    condensed here only to keep the complete interface visible on one page;
    \cref{fig:full-repl-contract} reproduces the exact initial prompt. \sid{S10} is the skill library:
    named procedures, each a \texttt{name}, a one-line \texttt{description} of when to
    consult it, and a \texttt{body} of ordered steps. Only the name/description index
    enters the system prompt, after \sid{S5}; a body is returned verbatim by
    \texttt{load\_skill(name)}, which the runner installs in the root and child REPLs
    when the list is non-empty, so a procedure costs tokens only on the turn that loads
    it. At most eight entries; \hzero{} carries none. \sid{S6}, \sid{S7} and \sid{S9}
    required a runtime injection point in the reference implementation of
    \citet{zhang2025rlm}; \sid{S10} adds one hook to record skill loads in the trace;
    the rest use its existing prompt and tool configuration. A proposal modifies
    exactly one surface.}
    \label{fig:edit-surfaces}
\end{figure}
    


\section{Self-Harness Optimization Pseudocode}
\label{app:algorithm}

\noindent
\fbox{%
\begin{minipage}{\dimexpr\linewidth-2\fboxsep-2\fboxrule\relax}
    \textbf{Algorithm: Self-Harness optimization} \\[0.3em]
    \textbf{Input:} harness $H_0 = \text{B1}$, held-in $\mathcal{D}_{\text{in}}$, held-out $\mathcal{D}_{\text{ho}}$, $T$, \textit{patience} \\[0.5em]
    \hrule
    \vspace{0.5em}
    \small
    \noindent $H \leftarrow H_0$; \enspace $\textit{rounds\_without\_promotion} \leftarrow 0$ \\
    \textbf{for} $t = 1 \dots T$ \textbf{do} \\
    \hspace*{1.5em} $\textit{traces} \leftarrow \text{run}(H, \mathcal{D}_{\text{in}}, \text{repetitions}=m)$ \hfill \textit{\% mining} \\
    \hspace*{1.5em} $\textit{failures} \leftarrow \text{sub\_verify}(\textit{traces})$; \enspace $\textit{patterns} \leftarrow \text{cluster}(\textit{failures})$ \\
    \hspace*{1.5em} $\textit{candidates} \leftarrow \text{propose}(H, \textit{patterns}, K)$ \hfill \textit{\% proposal} \\
    \hspace*{1.5em} $\textit{results} \leftarrow \left\{ \text{validate}(c, \mathcal{D}_{\text{in}} \cup \mathcal{D}_{\text{ho}}, \text{repetitions}=v) : c \in \textit{candidates} \cup \{H\} \right\}$ \\
    \hspace*{1.5em} $\textit{promoted} \leftarrow \left\{ c \in \textit{candidates} : \text{no\_regression}(\textit{results}[c], \textit{results}[H], \tau) \right.$ \\
    \hspace*{9.5em} $\left. \land\ \text{cost}(c) \in [\text{cost}_{\min}, \text{cost}_{\max}] \right\}$ \\
    \hspace*{1.5em} \textbf{if} $\textit{promoted} \neq \emptyset$ \textbf{then} \\
    \hspace*{3.0em} $H' \leftarrow \text{merge}(\textit{promoted})$ \\
    \hspace*{3.0em} $H \leftarrow H'$ \textbf{if} $\text{not regressed}(H', \textit{results})$ \textbf{else} $H$ \hfill \textit{\% merge failure promotes nothing} \\
    \hspace*{3.0em} $\textit{rounds\_without\_promotion} \leftarrow 0$ \\
    \hspace*{1.5em} \textbf{else} \\
    \hspace*{3.0em} $\textit{rounds\_without\_promotion} \mathrel{+}= 1$ \\
    \hspace*{3.0em} \textbf{if} $\textit{rounds\_without\_promotion} \ge \textit{patience}$ \textbf{then break} \\
    \textbf{return} $H$ \hfill \textit{\% frozen as SH-RLM}
\end{minipage}%
}

\section{Fine-tuned \rlm{} baseline (F1)}
\label{app:finetune}

F1 reproduces the weight-training arm of \citet{zhang2026harness} so that harness
optimization and weight optimization are compared on the same backbone, harness, and
environment. The recipe below follows their reported setup; any deviation forced by our
compute allocation will be recorded before training begins.

\begin{description}
  \item[Backbone.] Qwen3-30B-A3B-Instruct-2507, the model \citet{zhang2026harness}
    RL-train inside an \rlm{} harness, and the same backbone used for B1, $\lambda$-RLM, and SH-RLM so that F1
    differs from B1 only in its weights.
  \item[Algorithm.] RL with \texttt{prime-rl}: decoupled PPO with GRPO-style advantages
    and a KL penalty against the initial policy.
  \item[Reward.] The environment's deterministic verifier on the final answer, identical
    to the evaluation metric and to the promotion signal used by \sh{}, so the two
    optimizers are driven by the same outcome measure.
  \item[Rollouts.] Batch size $64$ with $4$ rollouts per sample.
  \item[Steps.] $150$ optimization steps for the length-generalization setting, with
    evaluation every $10$ steps; \citet{zhang2026harness} run $500$ steps, evaluating
    every $20$, for their cross-domain strategy experiments.
  \item[Training inputs.] Short instances only, at $8$k--$64$k tokens, drawn from the
    same source-environment short split used for harness optimization. Long instances are
    never trained on.
  \item[Evaluation.] The frozen checkpoint is evaluated on the same four untouched test
    sets as B1, $\lambda$-RLM, and SH-RLM, at $256$k--$2$M tokens on the long splits.
  \item[Hardware.] $8\times$$\lambda$-RLM00 nodes.
\end{description}

Two conditions govern whether F1 is reported at all. If the published checkpoint is
obtainable, we evaluate it directly rather than retraining, and say so. If neither the
checkpoint nor the compute for training is available, F1 is dropped and its absence is
reported alongside the published figures for the same environment and backbone. F1 is
budgeted and evaluated separately from the fixed-weight conditions (B1, $\lambda$-RLM, SH-RLM).

\section{Ablations}
\label{app:ablations}

\paragraph{Sub-verification.} Sub-verification enters the loop only through
weakness mining, so we can withhold it without changing anything else. We repeat
the full optimization run with the sub-verifier signal removed, giving the
proposer root outcomes and traces alone, and compare the resulting frozen harness
against SH-RLM on all four test sets. The comparison asks whether checkable
child-level evidence is what makes mined failure attributions actionable, or
whether the proposer recovers the same edits from traces without it. Because
sub-verification is post hoc, the same recorded runs can additionally be mined
under both conditions within a round, yielding paired evidence bundles that
differ only in the grounding signal; this trace-level comparison costs no extra
execution and complements the loop-level ablation. Because the full repeat
doubles optimization cost, it is a preregistered secondary run on the source
environment, contingent on the pilot's cost estimate holding.

\paragraph{Leave-one-edit-out.} Each promoted edit is removed from the final
harness individually and the affected test conditions are re-evaluated. These
ablations determine which edits are responsible for the final gains and whether
improvements arise from the accumulated harness rather than unrelated sampling
variation.

\section{Assessing Alignment Drift}
\label{app:alignment}
  
Given prior work on misevolution in self-optimizing agents~\citep{shao2025misevolution},
and in particular the finding that capability-driven workflow optimization can sharply
reduce refusal rates on harmful requests despite passing task-level
validation~\citep{shao2025misevolution}, an optional extension is to evaluate SH-RLM for
alignment drift as a side effect of harness promotion. Because our promotion gate, like
that of \citet{zhang2026selfharness}, is defined only over verifier accuracy and cost,
edits that are capability-positive but compliance-increasing (e.g., ``always commit to an
answer,'' suppressed hedging, or reasoning-mode changes, which are known to shift safety
behavior in the Qwen3 family~\citep{selfjailbreak2025}) would be promoted without
detection. The extension freezes a small safety probe of a few hundred harmful and
borderline prompts~\citep{wildjailbreak2024}, runs it through the \rlm{} under B1 and
under each promoted harness in the lineage, and scores refusal and harmful-answer rates
with an automatic judge, yielding a safety trajectory plotted alongside the accuracy
trajectory. At roughly $300$ prompts across $\sim$$8$ harness states, with probe runs
far shorter than benchmark instances, this adds a small fraction of the main evaluation
budget. It would enable us to assess
whether the propose--validate--accept structure incidentally preserves alignment, as its
bounded edit surfaces might suggest, or whether the harness pathway exhibits the same
drift documented for workflow evolution; either outcome is informative, and to our
knowledge neither has been measured for \sh{}-style harness optimization. Because no
promotion decision depends on the probe, it is purely observational and does not alter
the preregistered optimization loop.
\end{document}




  